\section{Discussion}

\subsection{Strengths}
The main strength of Examiner Coach is its modular backend design. The system
separates transcription, retrieval, evaluation, validation, and coaching, making
it easier to inspect and adapt individual components. The six-criterion rubric
also makes LLM output more structured and educationally interpretable.

Another strength is the direct connection between educational theory and system
design. The system does not merely ask an LLM to ``give feedback on feedback''.
It constrains the task through explicit criteria, score anchors, retrieved
evidence, and a structured output schema. This makes the output more suitable
for formative training because the examiner can see which aspect of the
feedback needs improvement and can ask targeted follow-up questions.

This directly addresses the training gap described in the problem definition.
Instead of relying only on occasional large-group examiner training, the
prototype supports repeated individual practice: an examiner can record a
two-minute feedback attempt, receive immediate criterion-based evaluation, and
then ask for clarification or alternative wording. The system therefore targets
the missing practice-and-correction loop rather than replacing formal examiner
training.

The modular design also supports experimentation. Retrieval can be disabled,
direct retrieval and HyDE can be compared, and filtering or reranking can be
switched on or off. This is valuable in a research project because it allows the
system to be studied as a set of computational modelling choices rather than as
a fixed black-box application.

\subsection{Limitations}
The current evaluation is prototype-level and should not be interpreted as a
final validation study. Results depend on transcription quality, retrieval
quality, prompt calibration, model availability, and the benchmark design. The
system does not replace expert human judgment.

The benchmark comparison is also limited by its size and construction. Expected
score bands provide a useful calibration target, but they are not equivalent to
consensus ratings from multiple medical education experts. The LLM-as-a-judge
analysis gives an additional consistency check, but it remains model-based and
therefore cannot replace human validation. A future study should compare system
outputs against trained human raters and examine inter-rater agreement.

Another limitation is model stability. The evaluation model used for the
notebook comparison is no longer available in the current KISSKI setup, so the
exact results should not be overinterpreted as final performance claims. The
architecture and comparison method remain valid, but the numerical results need
to be repeated under the current model configuration.

\subsection{Privacy and Ethical Considerations}
The system processes audio, transcripts, evaluation results, and coaching
interactions. External AI services may receive audio or text depending on the
configured endpoint. Future deployment requires clear retention policies, access
control, participant information, and review of institutional data protection
requirements.

This privacy aspect is part of the core problem, not an afterthought. In medical
education, examiner feedback may be patient-adjacent, institution-specific, or
linked to assessment material. For this reason, a tool intended for examiner
training cannot simply assume that commercial AI services are acceptable. The
prototype uses configured KISSKI/SAIA endpoints and a local knowledge base, but
future deployment would still need formal review of what data leaves the local
environment, how long audio and transcripts are retained, and who can access the
stored training outputs.

Ethically, the system should be positioned as a training aid rather than an
authority over examiner competence. Automated feedback evaluation can support
self-reflection, but it may also create false confidence if users treat scores
as objective judgments. For this reason, the report interprets scores as
formative indicators and emphasizes the need for expert review before live study
deployment.

\subsection{Future Work}
Future work should first validate the prototype with medical educators and OSCE
examiners. The current benchmark provides an initial calibration signal, but it
should be compared with human expert ratings and tested in a user study that
examines usefulness, trust, clarity of the criteria, and whether repeated
practice improves feedback delivery.

Technically, the next step is more detailed retrieval analysis. This includes
examining which evidence sources are retrieved, which rubric criteria benefit
most from RAG, and how retrieval behaves across German and English transcripts.
The comparison framework could also be extended to test different values of
\texttt{final\_k}, candidate-pool sizes, HyDE prompts, reranking rules, and
repeated model runs.

The extended comparison was not used in this report to avoid unnecessary use of
shared AI infrastructure and to remain mindful of the environmental impact of
large numbers of LLM and judge calls. Future experiments should therefore be
planned around specific research questions and reported together with the model,
number of calls, and configuration used.
